% Algorithm Comparison (PS3). Estilo visual de la presentacion PS2.
% Compilar con: source("04_slides/compile_slides.R") desde la raiz.
% Los numeros y graficos se generan en 01_code/14_compare_algorithms.R.
\documentclass[aspectratio=169,10pt]{beamer}
\usepackage{fontspec}
\setmainfont{Carlito}
\setsansfont{Carlito}
\usepackage[spanish,es-noshorthands]{babel}
\usepackage{graphicx,booktabs,array,tabularx,colortbl}
\usepackage{xcolor}
\definecolor{azulPS}{HTML}{2F5FA8}
\definecolor{tintaPS}{HTML}{293B44}
\definecolor{grisPS}{HTML}{61717B}
\definecolor{fondoPS}{HTML}{E4E7E9}
\setbeamercolor{normal text}{fg=tintaPS,bg=white}
\setbeamercolor{frametitle}{fg=white,bg=azulPS}
\setbeamercolor{structure}{fg=azulPS}
\setbeamerfont{frametitle}{family=\sffamily,series=\bfseries,size=\large}
\setbeamerfont{normal text}{size=\small}
\setbeamertemplate{navigation symbols}{}
\setbeamersize{text margin left=.65cm,text margin right=.65cm}
\setbeamertemplate{frametitle}{%
 \nointerlineskip
 \begin{beamercolorbox}[wd=\paperwidth,ht=.42cm,dp=.15cm,leftskip=.48cm]{frametitle}
  \usebeamerfont{frametitle}\insertframetitle
 \end{beamercolorbox}}
\setbeamertemplate{footline}{\hfill\color{tintaPS}\scriptsize\insertframenumber/14\hspace{.32cm}\vspace{.17cm}}
\setbeamertemplate{itemize item}{\raisebox{.08ex}{\tiny\textbullet}}
\setbeamertemplate{enumerate item}{\insertenumlabel.}
\newcommand{\fuente}[1]{\par{\fontsize{7}{8}\selectfont\color{grisPS}#1}}
\newcommand{\grafico}[2]{\includegraphics[width=\linewidth,height=#2,keepaspectratio]{../02_outputs/figures/#1.pdf}}
\newcommand{\nota}[1]{\par\medskip{\small #1}}
\newcommand{\cabeceraTabla}{\arrayrulecolor{tintaPS}\toprule}
\newcommand{\cierreTabla}{\bottomrule}
\setlength{\tabcolsep}{4pt}
\renewcommand{\arraystretch}{1.18}
\newcommand{\MAEBlend}{164,53}
\newcommand{\MAEVecinos}{166,73}
\newcommand{\MAEMdos}{169,16}
\newcommand{\MAELGBM}{178,61}
\newcommand{\MAERF}{174,74}
\newcommand{\MAECART}{208,28}
\newcommand{\MAERed}{180,01}
\newcommand{\MAEEN}{199,73}
\newcommand{\MAELM}{215,58}
\newcommand{\MAEBase}{296,62}
\newcommand{\MAETFIDFLGBM}{194,39}
\newcommand{\MAEDmin}{169,17}
\newcommand{\GainBlend}{2,20}
\newcommand{\GainTotal}{14,08}
\newcommand{\GainMdos}{9,45}
\newcommand{\GainVecinos}{2,43}
\newcommand{\GainMCO}{81,03}
\newcommand{\CVBlendAleatoria}{94,35}
\newcommand{\CVBlendEspacial}{128,58}
\newcommand{\CVGapEspacial}{9,80}
\newcommand{\CVGapAleatoria}{21,31}
\newcommand{\CVGapBloques}{9,88}
\newcommand{\NBase}{17}
\newcommand{\NCompleto}{174}
\newcommand{\NAmpliado}{12.174}

\begin{document}

% 1: Comienza con el mejor modelo, sin portada, segun el enunciado.
\begin{frame}[t]{Mejor modelo: BLEND de LightGBM y red neuronal}
\vspace{.20cm}
\begin{columns}[T,onlytextwidth]
\begin{column}{.48\textwidth}
\colorbox{fondoPS}{\parbox{.97\linewidth}{\vspace{.10cm}\bfseries Especificación\vspace{.10cm}}}
\vspace{.10cm}
\begin{itemize}\setlength\itemsep{.20cm}
\item 75\,\% LightGBM híbrido + 25\,\% red neuronal.
\item LightGBM combina precio por m$^2$, superficie y precios de 10 vecinos.
\item La red tiene tres capas ocultas: 256, 128 y 64 neuronas.
\item Peso elegido con CV espacial, sin vecinos.
\end{itemize}
\vspace{.07cm}
\begin{center}
{\fontsize{26}{27}\selectfont\color{azulPS}\bfseries \MAEBlend\,M}\par
{\small MAE en Kaggle (tabla pública)}\par
{\scriptsize CV espacial: \CVBlendEspacial\,M \quad CV aleatoria: \CVBlendAleatoria\,M}
\end{center}
\end{column}
\begin{column}{.49\textwidth}
\colorbox{fondoPS}{\parbox{.97\linewidth}{\vspace{.10cm}\bfseries Lo que muestra la comparación\vspace{.10cm}}}
\vspace{.12cm}
\begin{enumerate}\setlength\itemsep{.23cm}
\item \textbf{El modelo importa:} el BLEND supera a Random Forest (\MAERF\,M) y al LightGBM completo (\MAELGBM\,M) en Kaggle.
\item \textbf{Las transformaciones importan:} modelar precio por m$^2$ y añadir vecinos mejora al LightGBM completo.
\item \textbf{Combinar predicciones ayuda:} añadir la red reduce el MAE público del LightGBM con vecinos en \GainBlend\,M.
\end{enumerate}
\end{column}
\end{columns}
\vfill
\fuente{Fuente: kaggle\_scores.csv y registro\_modelos.csv. MAE en millones de COP.}
\end{frame}

% 2: Datos de la competencia.
\begin{frame}[t]{Datos: predecir precios en Chapinero sin observarlos en train}
\vspace{.38cm}
\begin{columns}[T,onlytextwidth]
\begin{column}{.49\textwidth}
\begin{tabular}{@{}lrr@{}}
\cabeceraTabla & \textbf{Train} & \textbf{Test} \\ \midrule
Propiedades & 38.644 & 10.286 \\
Zona & Resto Bogotá & Chapinero \\
Precio & Sí & No \\
\multicolumn{3}{l}{Objetivo: precio de venta (COP)}\\
\multicolumn{3}{l}{Métrica: MAE en pesos}\\
\cierreTabla
\end{tabular}
\vspace{.25cm}
\fuente{Competencia PS3 de MECA 4107; cifras de las bases del proyecto.}
\end{column}
\begin{column}{.47\textwidth}
\begin{itemize}\setlength\itemsep{.22cm}
\item No buscamos acertar en viviendas al azar: el test se concentra en Chapinero.
\item La información del resto de Bogotá ayuda, pero hay diferencias espaciales de precios y características.
\item El MAE penaliza errores en pesos: equivocarse en una vivienda cara cuesta más.
\item La pregunta es qué modelo predice mejor \emph{fuera de las zonas donde aprendió}.
\end{itemize}
\end{column}
\end{columns}
\end{frame}

% 3: Variables.
\begin{frame}[t]{Predictores: información de la vivienda, del entorno y del anuncio}
\vspace{.29cm}
\small
\begin{tabularx}{\textwidth}{@{}p{1.6cm}Xr@{}}
\cabeceraTabla \textbf{Conjunto} & \textbf{Qué agrega} & \textbf{Variables} \\ \midrule
Base & Superficie, cuartos, baños, tipo, ubicación, fecha e indicadores de calidad. & \NBase \\[.12cm]
Completo & + estrato por manzana, TransMilenio, parques, universidades, pendiente, distancias y señales extraídas del texto. & \NCompleto \\[.12cm]
Ampliado & + 12.000 términos TF-IDF (unigramas y bigramas de la descripción). & \NAmpliado \\
\cierreTabla
\end{tabularx}
\vspace{.46cm}
\colorbox{fondoPS}{\parbox{.96\textwidth}{\vspace{.12cm}\textbf{Hipótesis que podemos contrastar}\par\smallskip Más información puede reducir el error, pero no necesariamente: texto de alta dimensión añade ruido y cuesta más estimarlo. La comparación debe distinguir qué variables y qué algoritmo cambiaron.\vspace{.10cm}}}
\vfill\fuente{Fuente: conjuntos\_predictores.csv y scripts de preparación del repositorio PS3.}
\end{frame}

% 4: Tabla de algoritmos.
\begin{frame}[t]{Seis familias de modelos y el ensamble final}
\vspace{.25cm}
\begin{tabularx}{\textwidth}{@{}p{2.45cm}Xr@{}}
\cabeceraTabla \textbf{Familia} & \textbf{Variante con mejor envío público disponible} & \textbf{MAE Kaggle} \\ \midrule
Regresión lineal & MCO, conjunto completo & \MAELM\,M \\
Elastic Net & Texto TF-IDF, conjunto ampliado & \MAEEN\,M \\
CART & Árbol con precio por m$^2$ & \MAECART\,M \\
Random Forest & Bosque con precio por m$^2$ & \MAERF\,M \\
LightGBM & Precio por m$^2$ + vecinos & \MAEVecinos\,M \\
Red neuronal & Capas 256--128--64 & \MAERed\,M \\
\rowcolor{fondoPS}\textbf{BLEND} & \textbf{LightGBM con vecinos + red (75/25)} & \textbf{\MAEBlend\,M} \\
\cierreTabla
\end{tabularx}
\vspace{.28cm}
\nota{\textbf{Lectura:} el menor puntaje es del BLEND, pero las mejores variantes de cada familia no necesariamente usan la misma información ni la misma transformación del objetivo.}
\fuente{Fuente: Kaggle, puntaje público registrado en kaggle\_scores.csv.}
\end{frame}

% 5: Comparacion entre familias.
\begin{frame}[t]{Con el conjunto completo: Random Forest lidera la comparación inicial}
\vspace{.15cm}
\begin{columns}[T,onlytextwidth]
\begin{column}{.63\textwidth}
\grafico{compare_01_familias}{5.45cm}
\end{column}
\begin{column}{.33\textwidth}
\vspace{.25cm}
\begin{itemize}\setlength\itemsep{.20cm}
\item RF: \MAERF\,M, menor MAE público entre estas versiones.
\item LightGBM: \MAELGBM\,M; red: \MAERed\,M.
\item CART, Elastic Net y MCO obtienen errores mayores.
\end{itemize}
\vspace{.18cm}
\textcolor{azulPS}{\small\textbf{No basta decir ``boosting gana'':}} el liderazgo cambia cuando modificamos el objetivo y añadimos vecinos.
\end{column}
\end{columns}
\vfill
\fuente{Mismo conjunto completo de variables estructuradas. RF y CART modelan precio por m$^2$; otros, precio total. Comparación orientativa, no experimento puro de algoritmo.}
\end{frame}

% 6: Comparacion de features.
\begin{frame}[t]{Más variables ayudan, pero no siempre de la misma forma}
\vspace{.13cm}
\begin{columns}[T,onlytextwidth]
\begin{column}{.63\textwidth}
\grafico{compare_02_predictores}{5.48cm}
\end{column}
\begin{column}{.33\textwidth}
\vspace{.19cm}
\begin{itemize}\setlength\itemsep{.21cm}
\item En MCO, pasar de \NBase{} a \NCompleto{} variables reduce el MAE en \GainMCO\,M.
\item Elastic Net mejora con TF-IDF: 213,59 a \MAEEN\,M.
\item LightGBM empeora al añadir TF-IDF: \MAELGBM{} a \MAETFIDFLGBM\,M.
\end{itemize}
\smallskip
\textcolor{azulPS}{\textbf{Conclusión:}} más predictores no garantizan menor error.
\end{column}
\end{columns}
\vfill
\fuente{Fuente: puntajes públicos. Variantes de Elastic Net tienen ajustes distintos; diferencias no aíslan causalmente el efecto de TF-IDF.}
\end{frame}

% 7: Evolucion del modelo.
\begin{frame}[t]{LightGBM mejora al cambiar el objetivo y combinar predicciones}
\vspace{.25cm}
\begin{columns}[T,onlytextwidth]
\begin{column}{.67\textwidth}
\grafico{compare_03_evolucion}{5.30cm}
\end{column}
\begin{column}{.29\textwidth}
\begin{enumerate}\setlength\itemsep{.28cm}
\item Precio por m$^2$: cuando conocemos superficie real, predecimos precio unitario.
\item Vecinos: añadimos información de precios de otras propiedades.
\item BLEND: combinamos errores parcialmente distintos.
\end{enumerate}
\medskip
\textcolor{azulPS}{\textbf{Ganancia total observada: \GainTotal\,M}} frente al LightGBM completo.
\end{column}
\end{columns}
\vfill
\fuente{Fuente: Kaggle público. Secuencia de variantes desarrolladas; no experimento causal ni garantía de generalización.}
\end{frame}

% 8: Vecinos.
\begin{frame}[t]{A fondo: los vecinos ayudan, pero excluir cercanos no ayudó}
\vspace{.20cm}
\begin{columns}[T,onlytextwidth]
\begin{column}{.63\textwidth}
\grafico{compare_04_vecinos}{4.6cm}
\end{column}
\begin{column}{.33\textwidth}
\vspace{.10cm}
\begin{itemize}\setlength\itemsep{.24cm}
\item Agregar vecinos reduce el MAE de \MAEMdos{} a \MAEVecinos\,M.
\item Excluir referencias a menos de 200 m lo sube a \MAEDmin\,M.
\item La cercanía contiene señal valiosa; también exige vigilar anuncios duplicados y filtraciones.
\end{itemize}
\end{column}
\end{columns}
\vspace{.04cm}
\nota{\textbf{Precaución:} en CV, el precio de los vecinos debe provenir únicamente del subconjunto de entrenamiento del fold.}
\fuente{Fuente: tres envíos LightGBM por m$^2$ en Kaggle, script 11d y variante dmin = 200 m.}
\end{frame}

% 9: Peso de la red.
\begin{frame}[t]{A fondo: mezclar 25 \% de red reduce el MAE}
\vspace{.12cm}
\begin{columns}[T,onlytextwidth]
\begin{column}{.63\textwidth}
\grafico{compare_05_peso_red}{5.55cm}
\end{column}
\begin{column}{.33\textwidth}
\vspace{.20cm}
\begin{itemize}\setlength\itemsep{.26cm}
\item El peso 25\,\% minimiza el MAE en viviendas similares a Chapinero en la prueba \emph{sin vecinos}.
\item Ese peso se reutiliza en el BLEND con vecinos.
\item En Kaggle pasa de \MAEVecinos{} a \MAEBlend\,M: mejora de \GainBlend\,M.
\end{itemize}
\end{column}
\end{columns}
\vfill
\fuente{La curva corresponde a LGBM por m$^2$ + red sin vecinos (script 11e); el puntaje del BLEND final es el envío del 11f.}
\end{frame}

% 10: Validacion.
\begin{frame}[t]{Generalización: la validación aleatoria resulta optimista}
\vspace{.15cm}
\begin{columns}[T,onlytextwidth]
\begin{column}{.63\textwidth}
\grafico{compare_06_validacion}{5.45cm}
\end{column}
\begin{column}{.33\textwidth}
\vspace{.28cm}
\begin{itemize}\setlength\itemsep{.25cm}
\item En el BLEND, la CV aleatoria indica \CVBlendAleatoria\,M; la espacial, \CVBlendEspacial\,M.
\item Kaggle público registra \MAEBlend\,M.
\item El salto evidencia el reto de predecir un barrio diferente al de entrenamiento.
\end{itemize}
\end{column}
\end{columns}
\vfill
\fuente{Fuente: registro\_modelos.csv y kaggle\_scores.csv. MAE de CV puede ser optimista por selección de hiperparámetros y pesos.}
\end{frame}

% 11: Eleccion de CV espacial.
\begin{frame}[t]{Elegimos CV espacial simulando barrios parecidos a Chapinero}
\vspace{.25cm}
\begin{columns}[T,onlytextwidth]
\begin{column}{.64\textwidth}
\grafico{compare_07_cv_esquemas}{4.35cm}
\end{column}
\begin{column}{.32\textwidth}
\textbf{Criterio de elección}
\begin{itemize}\setlength\itemsep{.15cm}
\item Creamos pseudo-Chapineros dentro de train.
\item Comparamos error prometido por CV y error real al reservar esas zonas.
\item K-means de 20 regiones obtiene menor brecha media absoluta.
\end{itemize}
\end{column}
\end{columns}
\nota{\textbf{Implicación:} una buena CV no es la que da el menor MAE, sino la que mejor anticipa una zona nueva. La ventaja frente a bloques de 700 m fue estrecha: \CVGapEspacial\,\% vs. \CVGapBloques\,\%.}
\fuente{Fuente: seleccion\_cv.csv; evaluación con Random Forest y MCO en pseudo-Chapineros.}
\end{frame}

% 12: Tradeoffs.
\begin{frame}[t]{¿Qué modelo elegiríamos y qué sacrificamos?}
\vspace{.26cm}
\scriptsize
\begin{tabularx}{\textwidth}{@{}p{2.25cm}p{1.15cm}X X@{}}
\cabeceraTabla \textbf{Modelo} & \textbf{MAE público} & \textbf{Ventaja} & \textbf{Costo o límite} \\ \midrule
\rowcolor{fondoPS}\textbf{BLEND (75/25)} & \textbf{\MAEBlend\,M} & Menor error público & Dos modelos, menor interpretabilidad \\
LightGBM + vecinos & \MAEVecinos\,M & Alta precisión & Usa precios de inmuebles cercanos \\
Random Forest & \MAERF\,M & Flexible; un solo modelo & Interpretación compleja \\
CART & \MAECART\,M & Reglas legibles & Error mucho mayor \\
MCO completo & \MAELM\,M & Coeficientes simples & Menos flexibilidad \\
\cierreTabla
\end{tabularx}
\vspace{.39cm}
\begin{columns}[T,onlytextwidth]
\begin{column}{.48\textwidth}
\normalsize \textbf{Competencia:} elegiríamos el BLEND por su menor MAE público.
\end{column}
\begin{column}{.48\textwidth}
\normalsize \textbf{Uso inmobiliario:} también importan estabilidad, explicabilidad y sobreestimaciones.
\end{column}
\end{columns}
\vfill
\fuente{Fuente: Kaggle público. No se atribuyen tiempos computacionales no medidos de forma comparable.}
\end{frame}

% 13: Mensajes principales.
\begin{frame}[t]{Tres conclusiones de la comparación}
\vspace{.43cm}
\begin{enumerate}\setlength\itemsep{.36cm}
\item No ganamos solo por escoger un algoritmo: la \textbf{ingeniería de variables}, el cambio de objetivo a precio por m$^2$ y el ensamble explican buena parte de las mejoras observadas.
\item La evaluación debe representar el problema real: \textbf{predecir Chapinero} es más difícil que validar aleatoriamente propiedades de Bogotá.
\item El \textbf{BLEND gana en el leaderboard público}: \MAEBlend\,M de MAE. Para decisiones de compra necesitamos contrastar el resultado con validación externa y controlar la sobreestimación.
\end{enumerate}
\vspace{.30cm}
\begin{center}\color{azulPS}\large\bfseries El mejor puntaje es el punto de partida, no el final de la evaluación.\end{center}
\vfill\fuente{Lectura de resultados versionados del PS3; el leaderboard privado puede diferir del público.}
\end{frame}

% 14: Despedida.
\begin{frame}[plain]
\begin{center}
\vspace{2.0cm}
{\fontsize{34}{40}\selectfont\color{azulPS}\bfseries ¡Gracias!}\par\vspace{.43cm}
{\large Equipo 03: Julián Herrera, Andrés Silva y Valentina Vera}\par\vspace{.15cm}
{\small MECA 4107 · Algorithm Comparison · 2026-20}
\end{center}
\end{frame}
\end{document}
